\src{184}
\result{LEMMA}{5} Let $A$ and $B$ be two coherent algebraic sheaves on $X$, and
set $A(m)=A\otimes\Osh_X(m)$. Then, for $m$ sufficiently large, the canonical
homomorphism
\[
  \operatorname{Ext}_{\Osh_X}^{1}(X;A(-m),B)
  \longrightarrow H^{0}\bigl(X,\Extsh_{\Osh_X}^{1}(A(-m),B)\bigr)
  =H^{0}\bigl(X,\Extsh_{\Osh_X}^{1}(A,B)(m)\bigr)
\]
is an isomorphism.

This follows at once from the spectral sequence of Proposition 1 applied to
$A(-m)$ and $B$, since one then has
\[
  E_{2}^{p,q}(A(-m),B)
  =H^{p}\bigl(X,\Extsh_{\Osh_X}^{q}(A(-m),B)\bigr)
  =H^{p}\bigl(X,\Extsh_{\Osh_X}^{q}(A,B)(m)\bigr),
\]
which vanishes for $p>0$ and $m$ large.

We now prove (D) in the case $X=P$. First we prove that (5.2 bis) is an
isomorphism for $p=n$. Both sides are then left-exact functors (since
$H^{r+1}(P,F)=0$), so Lemma 3 shows that it suffices to prove the assertion for
$F=\Osh_P(-m)$; in that case it is contained in Lemma 2. The homomorphisms
(5.2 bis) are functorial and compatible with the connecting homomorphisms, and
for $p<n$ its two sides are coeffaceable functors of $F$ (Corollary to Lemma 4).
A standard argument therefore shows that (5.2 bis) is an isomorphism for every
$p$. This proves the duality theorem for projective space. Now let $X$ be
arbitrary but nonsingular. By the duality theorem for $P$, there is an
isomorphism
\[
  H^{n}(X,F)'=H^{n}(P,F)'
  =\operatorname{Ext}_{\Osh_P}^{r-n}(P;F,\Omega_P^{r}).
\]
By Lemma 1 (no. 4), the last term identifies with
\[
  \operatorname{Hom}_{\Osh_P}(P;F,\omega_X^{n})
  =\operatorname{Hom}_{\Osh_X}(X;F,\Omega_X^{n})
  =\operatorname{Ext}_{\Osh_X}^{0}(X;F,\Omega_X^{n}),
\]
and hence there is an isomorphism
\begin{equation}\tag{5.6}
  H^{n}(X,F)'=\operatorname{Hom}_{\Osh_X}(X;F,\Omega_X^{n})
  =\operatorname{Ext}_{\Osh_X}^{0}(X;F,\Omega_X^{n}).
\end{equation}
Taking $F=\Omega_X^{n}$ gives an isomorphism
\begin{equation}\tag{5.7}
  \gamma:H^{n}(X,\Omega_X^{n})\xrightarrow{\sim}k.
\end{equation}

